\documentclass[11pt,a4paper]{article}
\usepackage[T1]{fontenc}
\usepackage[utf8]{inputenc}
\usepackage[french]{babel}
\frenchsetup{SmallCapsFigTabCaptions=false}
\usepackage{lmodern,microtype,amsmath}
\usepackage[margin=23mm]{geometry}
\usepackage{booktabs,array,tabularx,longtable}
\usepackage{xcolor,tikz,enumitem,float}
\usepackage{hyperref}
\usetikzlibrary{arrows.meta,positioning,fit,calc}
\definecolor{ink}{HTML}{1F242C}
\definecolor{muted}{HTML}{5C626C}
\definecolor{line}{HTML}{D9DEE5}
\definecolor{accent}{HTML}{AF0032}
\definecolor{rose}{HTML}{FDF0F4}
\hypersetup{colorlinks=true,linkcolor=black,citecolor=black,urlcolor=ink,
 pdftitle={LANCE : architecture, protocoles et état d'avancement},
 pdfauthor={Projet LANCE},pdflang={fr-FR}}
\urlstyle{same}
\def\UrlBreaks{\do\/\do\-\do\_\do\.\do\:}
\newcolumntype{Y}{>{\raggedright\arraybackslash}X}
\DeclareRobustCommand{\code}[1]{{\small\path{#1}}}
\newcommand{\source}[2]{\href{../../#1}{\nolinkurl{#1}}. #2}
\newcommand{\revision}{2026-10-30 09:00 (Europe/Brussels, UTC+01:00)}
\setlength{\parskip}{3pt}
\setlength{\emergencystretch}{2em}
\setlist{nosep,leftmargin=5mm,topsep=3pt}
\clubpenalty=10000 \widowpenalty=10000
\setcounter{tocdepth}{1}

\title{LANCE\\[4pt]
\large Architecture, protocoles et état d'avancement\\[2pt]
\normalsize Rapport au 30 octobre 2026}
\author{Projet LANCE\\\small Rédaction par volets spécialisés, relecture et synthèse indépendantes}
\date{Dernière révision : \revision\\[3pt]
\small État documentaire au 30 octobre 2026 ; aucune expérience nouvelle rapportée}

\begin{document}
\maketitle

\begin{abstract}
LANCE est une plateforme d'audit de sécurité de réseaux d'objets connectés
autorisés, développée dans le cadre des travaux UMONS du programme NATO SPS
G6259. Ce rapport décrit l'architecture du pipeline, les protocoles d'audit et
leurs contrats de preuve, les harnais d'expérimentation et leur bac à sable
hermétique, l'infrastructure d'exécution et le circuit de livraison, puis
l'état réel du benchmark. Il succède au rapport d'ensemble du 15 octobre 2026,
qui survolait l'architecture et les protocoles : les parties correspondantes
sont ici développées jusqu'au contrat de code.

Le document distingue systématiquement ce qui est implémenté, ce qui est
mesuré et ce qui reste proposé. Aucune expérience nouvelle n'est rapportée :
les mesures citées sont celles déjà consignées, avec leur date. Les limites de
chaque composant sont explicites, y compris celles du bac à sable, dont le
garde-fou n'est pas un bac à sable natif exhaustif.
\end{abstract}

\tableofcontents
\clearpage

\input{partie-intro}

\section{Objet et méthode}

Ce rapport répond à une question de lecture : \emph{comment le système est-il
construit, et où en est-il réellement ?} Il ne remplace pas les pages de
référence de \code{docs/}, qui décrivent le comportement courant, ni les
protocoles détaillés de \code{research/}, qui portent le raisonnement.

Chaque partie a été rédigée par un rédacteur spécialisé travaillant sur le code
de la branche \code{dev/1}, puis relue contre ce même code par un relecteur
indépendant. Une synthèse transverse a été produite séparément. Les affirmations
non vérifiables dans le dépôt sont signalées comme telles plutôt que comblées.

Le rapport est daté : il décrit un état, pas un contrat. Un comportement
ultérieur peut le contredire sans que le rapport soit faux.

% === PARTIES RÉDIGÉES PAR LES AGENTS ===
\input{partie-architecture}
\input{partie-protocoles}
\input{partie-harnais}
\input{partie-infrastructure}
\input{partie-benchmark}

\section{Conclusion}

% === SYNTHÈSE ===
\input{partie-synthese}

\end{document}
